\documentclass[11pt]{article}
\usepackage[margin=1in]{geometry}
\usepackage{amsmath,amssymb,amsthm}
\usepackage[colorlinks=true,urlcolor=blue,linkcolor=blue,citecolor=blue]{hyperref}
\newtheorem{theorem}{Theorem}
\newtheorem{lemma}[theorem]{Lemma}
\newtheorem{corollary}[theorem]{Corollary}
\newcommand{\R}{\mathbb R}
\newcommand{\norm}[1]{\left\lVert#1\right\rVert}
\newcommand{\rank}{\operatorname{rank}}
\newcommand{\range}{\operatorname{range}}
\newcommand{\tO}{\widetilde O}
\newcommand{\MM}{T_{\rm MM}}
\newcommand{\Ac}{\mathcal A}
\title{Block spectral approximation of a factored product\\
from regularized absolute-value pencils}
\author{Working manuscript: independently checked proof; priority under review}
\date{November 14, 2026}
\begin{document}
\maketitle
\begin{abstract}
For arbitrary real factors \(A\in\R^{n\times m}\) and
\(B\in\R^{m\times p}\), we construct a width-at-most-\(r\) approximation
with spectral error at most \((1+\theta)\sigma_{r+1}(AB)\), with probability
at least \(1-\delta\). A coarse Gram sketch of relative accuracy \(h\)
preconditions every positive spectral threshold, including singular ones.
The finite preconditioner has a weighted quadratic defect of order \(h\).
A strong block residual certificate transfers to the original product
with error of order \(\alpha h\), rather than \(\alpha\). This permits
\(\widetilde O(\sqrt{h/\theta})\) true-factor block actions.
If the raw Krylov width is at most \(m\), total arithmetic cost is
\[
 \tO\!\left(I+G_h+m^{\omega+\zeta}
       +\sqrt{h/\theta}\,[\MM(n,m,r+1)+\MM(p,m,r+1)]
       +np+\MM(n,r,p)\right).
\]
Here \(I=m(n+p)\) and \(G_h\) is the cost of two sketched Grams.
Every output entry is included. An explicit regime permits every target
rank up to the square root of the outer dimension with exponent-two total
output cost, improving the checked fast batched Krylov and exact-Gram bounds.
A fixed-threshold variant targets absolute error without a supplied rank.
The result counts exact real arithmetic, scalar square roots, comparisons,
and Gaussian samples; it makes no finite-precision claim. Absolute-value
preconditioning, inertia, polynomial filters, and block Krylov are established
tools. Global priority of the combined theorem remains under review.
\end{abstract}

\section{Theorem and arithmetic model}
Let \(M=AB\), \(m\leq\min(n,p)\), \(1\leq r\leq m\), and
\(I=m(n+p)\). Write \(\MM(a,b,c)\) for an available arithmetic cost
of multiplying an \(a\)-by-\(b\) matrix by a \(b\)-by-\(c\) matrix.
We choose these costs symmetrically under permutation of the three
dimensions, as permitted by the matrix multiplication tensor.
Eigenvalues and singular values past a matrix's dimension are taken to be zero.
Put
\begin{align*}
 \Ac(s)&=\MM(n,m,s)+\MM(p,m,s),\\
 s_h&=\tO(m/h^2),\\
 G_h&=\MM(m,\min(n,s_h),m)+\MM(m,\min(p,s_h),m).
\end{align*}
The sketch dimensions include the stated failure probability. If a sketch
would have more rows than its input, use the identity instead.
We use any feasible square multiplication exponent \(\omega>2\), with
arbitrary fixed slack \(\zeta>0\); in particular \(\omega\leq9/4\)
is available over the reals \cite{amanov}.
Tildes hide only logarithms in dimensions, \(1/\theta\), and \(1/\delta\).
There is no hidden singular-value-ratio or factor-condition-number logarithm.
Field operations, exact comparisons, positive scalar square roots, and
Gaussian sampling have unit cost. The sizes of their representations are
not bounded. This distinction matters for characteristic coefficients,
exact sign tests, and the polynomial recurrences in the proof.

\begin{theorem}[Block product approximation]\label{thm:main}
Fix \(0<\theta\leq h\leq1/64\) and \(0<\delta<1\).
An algorithm returns factors of width at most \(r\) and explicitly forms
their product \(C\), such that, for every fixed pair \(A,B\),
\[
 \Pr\{\norm{M-C}_2\leq(1+\theta)\sigma_{r+1}(M)\}\geq1-\delta.
\]
Let \(q=\tO(\sqrt{h/\theta})\) be a sufficient capped per-run degree
for the block Krylov and final filter below, and put
\(L=(2q+1)(r+1)\). If \(L\leq m\), its total cost is
\begin{equation}\label{eq:main}
 \tO\!\left(I+G_h+m^{\omega+\zeta}+q\Ac(r+1)
                          +np+\MM(n,r,p)\right).
\end{equation}
For arbitrary \(L\), the same algorithm has the conservative bound
\begin{equation}\label{eq:general}
 \tO\!\left(I+G_h+m^{\omega+\zeta}+q\Ac(r+1)+mL^2
                          +np+\MM(n,r,p)\right).
\end{equation}
The condition uses the actual chosen degree, including its logarithms.
No rank, conditioning, gap, positive-tail, or factor-energy promise is required.
\end{theorem}

The proof constructs a physical subspace and returns its orthogonal
projection of \(M\). It uses explicit sketched Grams and true-factor
actions, without an exact spectral decomposition. If the product rank is at most \(r\),
it recovers \(M\) exactly on the common sketch event. Iteration counts
and dimensions are capped in advance; the arithmetic bounds also hold
on failed random outcomes.

\section{Sketches, support, and exact range algebra}
Set \(P=A^\top A\) and \(Q_{\rm raw}=BB^\top\). Fast subspace embeddings
\cite{cnw} produce explicit sketched Grams \(\bar P,\bar Q\) with
\begin{equation}\label{eq:ose}
 (1-h)P\preceq\bar P\preceq(1+h)P,\qquad
 (1-h)Q_{\rm raw}\preceq\bar Q\preceq(1+h)Q_{\rm raw}
\end{equation}
on an event with failure probability at most \(\delta/2\), in
\(\tO(I+G_h)\) operations. For example, dense fast transforms and
row sampling with \(\tO(m/h^2)\) rows suffice. Padding and an identity
choice handle all ambient dimensions.

Let \(a=\rank\bar P\). Exact rank revelation followed by Cholesky
computes \(W_0\in\R^{m\times a}\) and \(L_0=W_0^\dagger\) with
\(\range W_0=\range\bar P\) and \(W_0^\top\bar P W_0=I\).
Compute an orthonormal range basis \(U_0\) of
\(L_0\bar QL_0^\top\), and define
\begin{gather*}
 \widehat k=\operatorname{width}(U_0),\quad
 W=W_0U_0,\quad L_c=U_0^\top L_0,\quad Z=AW,\\
 G=Z^\top Z,\quad Q=L_cBB^\top L_c^\top,\quad Q_0=L_c\bar QL_c^\top.
\end{gather*}
Only \(W,L_c,Q_0\) are materialized; \(Z,G,Q\) are operators.
All this core algebra costs \(\tO(m^{\omega+\zeta})\).
On \eqref{eq:ose}, \(\widehat k=k=\rank M\) and
\begin{equation}\label{eq:support}
 M=ZL_cB,\qquad
 \frac{I}{1+h}\preceq G\preceq\frac{I}{1-h},\qquad
 (1-h)Q\preceq Q_0\preceq(1+h)Q.
\end{equation}
The identity for \(M\) follows because \(A=AW_0L_0\), and the columns
of \(L_0B\) lie in \(\range U_0\). In general \(A=AWL_c\) is false;
only the product identity is used. The nonzero spectrum of \(MM^\top\)
is that of \(H=G^{1/2}QG^{1/2}\).

The dimension cap is valid even off the sketch event. A sketched Gram
always has \(\range W_0\subseteq\range P\), and \(A\) is injective
on this subspace. If \(S_B\) denotes the right-factor embedding, then
\[
 \widehat k=\rank(L_0B S_B^\top)
       \leq\rank(\Pi_{\range P}B)=\rank(AB).
\]
Also \(Z\) has independent columns and \(Q_0\) is positive definite
on the computed support on every outcome. If \(\widehat k=0\), write
zero. If \(\widehat k\leq r\), form an orthonormal range basis of \(Z\)
and project \(M\) onto it. This costs \(\tO(\Ac(r)+m^{\omega+\zeta})\)
before expansion and is exact on the sketch event. Henceforth the good-event
analysis has \(r<k=\widehat k\), so \(\sigma_{r+1}(M)>0\).

For clarity, all orthogonalization steps have finite algebraic algorithms.
The deterministic CUP algorithm \cite{jps}, applied to a transpose,
selects a maximal independent column subset of any matrix, with no
generic-rank assumption. For a full-column-rank \(a\)-by-\(s\) matrix
\(V\), form \(V^\top V=R^\top R\) and return \(VR^{-1}\).
Fast recursive Cholesky and triangular inversion cost
\(O(s^{\omega+\zeta})\); their recurrences use block multiplication
and positive scalar square roots. Total range-basis cost is
\begin{equation}\label{eq:qr}
 \MM(s,a,s)+O(s^{\omega+\zeta})+\MM(a,s,s),
\end{equation}
in addition to rank revelation. In particular all dimensions at most
\(m\) cost \(\tO(m^{\omega+\zeta})\). No stability assertion about
this exact Gram construction is intended.

\section{One fixed inverse mass and a finite absolute preconditioner}
Put
\[
 \varepsilon_m=\theta/10^6,\qquad \alpha=\theta/(10^6h).
\]
Thus \(0<\alpha\leq10^{-6}\). Choose an even integer \(l\) with
\((h/(1-h))^{l+1}\leq\varepsilon_m\), and define the fixed operator
\begin{equation}\label{eq:mass}
 B_m=\sum_{i=0}^{l}(I-G)^i.
\end{equation}
Its degree is \(O(\log(1/\theta))\). Since
\(B_mG=I-(I-G)^{l+1}\), \eqref{eq:support} gives
\begin{equation}\label{eq:massbounds}
 (1-\varepsilon_m)G^{-1}\preceq B_m\preceq(1+\varepsilon_m)G^{-1},
 \qquad (1-2h)I\preceq B_m\preceq(1+2h)I.
\end{equation}
It is important to use this same polynomial at every threshold.
Even off the sketch event, it is positive definite: for any \(g>0\)
and even \(l\), its scalar value is
\([1-(1-g)^{l+1}]/g>0\).
The analytical matrix \(\widehat H=B_m^{-1/2}QB_m^{-1/2}\) satisfies
\begin{equation}\label{eq:hatcompare}
 \frac{\lambda_j(H)}{1+\varepsilon_m}
 \leq\lambda_j(\widehat H)
 \leq\frac{\lambda_j(H)}{1-\varepsilon_m}.
\end{equation}
No matrix root in this definition is computed. For every \(t>0\),
the pencil \(K_t=tB_m-Q\) has exactly as many negative eigenvalues
as \(\widehat H\) has eigenvalues greater than \(t\).

\begin{lemma}[Finite weighted defect]\label{lem:defect}
At every positive \(t\), one can form a positive definite core
\(F=JJ^\top\), together with \(J^{-1}\), in
\(\tO(k^{\omega+\zeta})\) arithmetic operations. On the sketch event,
with \(D=tI+Q_0\), it satisfies
\begin{equation}\label{eq:defect}
 8hD\preceq F\preceq2D,\qquad -F\prec K_t\prec F,\qquad
 0\preceq\Delta:=F-K_tF^{-1}K_t\preceq36hD.
\end{equation}
Moreover \(F-K_t\succeq6hD\).
\end{lemma}
\begin{proof}
Materialize \(D\) and the commuting symmetric matrix
\(S=(tI-Q_0)D^{-1}\), whose spectrum lies in \([-1,1]\).
Equation \eqref{eq:massbounds} and the proxy bounds imply
\[
 D^{-1/2}K_tD^{-1/2}=S+E,\qquad \norm E_2\leq2h.
\]
Indeed, the two errors are bounded by \(2htI\) and \(2hQ_0\)
before normalization. This bound has no \(\lambda_1/t\) factor.
Let \(a=8\) and \(M_0=S^2+a^2h^2I\). Starting at \(X_0=2I\),
use exact Newton steps
\[
 X_{i+1}=\tfrac12(X_i+M_0X_i^{-1}).
\]
These matrices commute with \(S,D\). For an eigenvalue
\(\rho=\sqrt{\lambda(M_0)}\leq2\), a scalar iterate \(x\geq\rho\)
obeys \(x_{\rm next}-\rho=(x-\rho)^2/(2x)\leq(x-\rho)/2\).
A predetermined \(O(\log(1/h))\) steps therefore give
\(\sqrt{M_0}\preceq X\preceq\sqrt{M_0}+ahI\).
Set \(F=D(X+ahI)\), compute its Cholesky factor, and invert that factor.
These are finite operations, without a matrix-square-root oracle.

In normalized coordinates, \(\bar F=X+ahI\) lies between
\(F_0=\sqrt{S^2+a^2h^2I}+ahI\) and \(F_0+ahI\).
Thus \(8hI\preceq\bar F\preceq2I\), and its margin over \(|S|\)
is at least \(8hI\). The perturbation \(E\) gives
\(-\bar F\prec S+E\prec\bar F\), with margin at least \(6hI\).
For the defect, the commuting scalar identity is
\(F_0-SF_0^{-1}S=2ahI\). The derivative of \(x-s^2/x\) is at most
two for \(x\geq|s|\), so
\(\bar F-S\bar F^{-1}S\preceq4ahI\).
In expanding \((S+E)\bar F^{-1}(S+E)\), the two cross terms cost
at most \(4hI\), because \(\norm{S\bar F^{-1}}_2\leq1\);
the remaining negative \(-E\bar F^{-1}E\) can be dropped.
This proves the upper bound \((4a+4)hI=36hI\).
The lower defect bound follows by congruence from
\(-\bar F\prec S+E\prec\bar F\). Undoing normalization proves the lemma.
\end{proof}

Define the actual symmetric operators
\begin{equation}\label{eq:CT}
 C_t=J^{-1}K_tJ^{-\top},\qquad T_t=I-C_t.
\end{equation}
On the sketch event, \(-I\prec C_t\prec I\) and
\(3hI\preceq T_t\preceq2I\), by \eqref{eq:defect}.
All block actions use \eqref{eq:mass}, the original factors, and explicit
cores, costing \(\tO(\Ac(b))\) on a block of width \(b\).
The inverses of the explicit triangular factors were included in setup.
This construction is motivated by established absolute-value preconditioning
\cite{vk,vy,plhr}. The weighted defect, rather than a spectral gap around
zero, is the certificate used here.

\section{Threshold decisions from exact coefficients}
Only comparisons with two fixed spectral levels are needed at a query.
Characteristic coefficients provide these comparisons and a positive
initial bracket, without computing ordered eigenvalues.

\begin{lemma}[Coefficient ratio and shifted inertia]\label{lem:coeff}
For a positive semidefinite \(d\)-by-\(d\) matrix \(V\), let
\(v_1\geq\cdots\geq v_d\geq0\) be its eigenvalues, and let \(e_i\)
be its elementary symmetric coefficients. If \(v_j>0\), then
\begin{equation}\label{eq:coeff}
 v_j/j\leq e_j/e_{j-1}\leq(d-j+1)v_j.
\end{equation}
The ratio can be computed in \(\tO(d^{\omega+\zeta})\) operations.
For any real symmetric matrix \(E\) and real \(s\), the number of
eigenvalues strictly greater than \(s\) can also be computed in that
cost, without eigenvalue separation or multiplicity assumptions.
\end{lemma}
\begin{proof}
The identity
\(j e_j=\sum_{|I|=j-1}v_I\sum_{i\notin I}v_i\) gives the lower
bound, since every \(I\) misses at least one of the top \(j\) indices.
For the upper bound, remove the largest-index element from each
size-\(j\) subset. Its value is at most \(v_j\), and each remaining
size-\(j-1\) subset is charged at most \(d-j+1\) times.
Compute the characteristic polynomial by a deterministic fast algorithm
over the reals \cite{np,npv}; its coefficients give \(e_i\) exactly.
If \(e_j=0\), then \(v_j=0\), which is checked before dividing.

For inertia, form \(p(x)=\det(xI-(E-sI))\). All its roots are real.
Let \(V(p)\) denote coefficient sign variations with zeros omitted.
Descartes' rule is exact in this case, including multiplicities.
For completeness, divide out any zero roots. If the remaining degree is
\(D\) and the positive and negative root counts are \(P,N\), Descartes
gives \(P\leq V(p)\), \(N\leq V(p(-x))\), while
\(V(p)+V(p(-x))\leq D\). To see the latter inequality, each gap of
\(g\) between consecutive nonzero coefficient degrees contributes at
most \(g\) to the sum: one when \(g\) is odd, and zero or two when
\(g\) is even. Since \(P+N=D\), all bounds are equalities.
Thus \(V(p)\) counts exactly the eigenvalues of \(E\) greater than \(s\).
It needs only the fast characteristic polynomial and \(O(d)\) sign tests.
An eigenvalue exactly equal to \(s\) gives a zero root and is excluded,
as required. No root isolation or spectral-decomposition oracle is used.
\end{proof}

Set \(j=r+1\) and \(\lambda_*=\lambda_j(\widehat H)>0\).
By Lemma~\ref{lem:coeff}, the plain coefficient ratio
\[
 \rho=\frac{e_j(Q_0)}{e_{j-1}(Q_0)}
 \quad\hbox{satisfies}\quad
 \rho/k\leq\lambda_j(Q_0)\leq k\rho.
\]
Moreover, the generalized eigenvalue principle gives
\[
 \frac{\lambda_j(Q_0)}{(1+h)(1+2h)}
 \leq\lambda_*\leq
 \frac{\lambda_j(Q_0)}{(1-h)(1-2h)}.
\]
Thus \([\rho/(4k),4k\rho]\) brackets \(\lambda_*\) with ratio \(16k^2\).
This initialization costs \(\tO(k^{\omega+\zeta})\) and does not
search from the largest eigenvalue to the target scale. A polynomial
dimension ratio is sufficient: bisection needs only dimension logarithms.

\begin{lemma}[Safe threshold oracle]\label{lem:quantile}
For each fixed positive \(t\), a block procedure returns one of
\emph{below}, \emph{above}, or \emph{near}, respectively certifying
\begin{equation}\label{eq:quantile}
 t<\lambda_*,\qquad t\geq\lambda_*,\qquad |\mu_j(C_t)|\leq\alpha,
\end{equation}
except with probability \(\beta\), conditional on the sketch event.
Here \(\mu_i\) orders eigenvalues increasingly. A \emph{below} decision
is correct on every run on that event. The procedure
uses \(\tO(\alpha^{-1/2})\) width-\(j\) true-factor block actions,
plus exact range algebra and \(\tO(m^{\omega+\zeta})\) core work.
\end{lemma}
\begin{proof}
For \(T=T_t\), form the Musco--Musco block Krylov range
\[
 \range V=\operatorname{span}[T\Omega,T^3\Omega,\ldots,T^{2q+1}\Omega],
 \qquad \Omega\in\R^{k\times j}\ \hbox{Gaussian},
\]
and take any exact orthonormal basis \(V\), of width \(d\leq\min(k,L)\).
Set \(Y=TV\), \(A_s=Y^\top Y\), and \(\varepsilon=\alpha/100\).
The per-vector guarantee, Theorem~12 of \cite{mm}, implies with
probability \(99/100\) that
\begin{equation}\label{eq:pervector}
 \beta_j:=\lambda_j(A_s)
 \geq\lambda_j(T)^2-\varepsilon\lambda_{j+1}(T)^2
 \geq(1-\varepsilon)\lambda_j(T)^2.
\end{equation}
This is an application of the source's per-vector theorem, not merely
its spectral low-rank guarantee. The compressed singular values are
the same as in its final Ritz extraction, although we do not compute
the Ritz vectors. Deterministically \(\beta_j\leq\lambda_j(T)^2\).

Repeat independently \(O(\log(1/\beta))\) times. Conceptually let
\(\beta_{\max}\) be the maximum resulting \(\beta_j\); define a missing
\(j\)-th eigenvalue as zero. It is unnecessary to compute this value.
Lemma~\ref{lem:coeff} tests only whether each \(\beta_j>1\) and whether
each \(\beta_j\leq s_0=1-2\varepsilon\), using exact shifted inertia.
If any \(\beta_j>1\), return \emph{below}. This is deterministically
safe because \(\lambda_j(T)^2\geq\beta_j>1\).
If all \(\beta_j\leq s_0\), return \emph{above}. When any repetition
succeeds in \eqref{eq:pervector},
\(\lambda_j(T)^2\leq s_0/(1-\varepsilon)<1\), so this is safe.
Otherwise \(1-2\varepsilon<\beta_{\max}\leq1\); with any successful
repetition,
\[
 1-2\varepsilon<\lambda_j(T)^2\leq(1-\varepsilon)^{-1}.
\]
Thus \(|1-\lambda_j(T)|\leq2\varepsilon<\alpha\), giving \emph{near}.
Here \(\mu_j(C_t)=1-\lambda_j(T)\). All repetitions fail with probability
at most \((1/100)^{O(\log(1/\beta))}\), proving the stated amplification.
Each repetition uses a separate basis; no eigenvalue validator, numerical
Ritz extraction, matrix powering, or root isolation is needed. Exact
characteristic coefficients and their sign tests cost
\(\tO(d^{\omega+\zeta})\). The range-basis cost is charged below.
\end{proof}

\section{A fresh projection with a strong residual}
A separate final filter avoids any assumption that truncating a larger
Krylov space preserves a previously obtained operator-norm certificate.

\begin{lemma}[Filtered Gaussian projection]\label{lem:filter}
Suppose \(0\preceq T\preceq2I\) and
\(\lambda_{r+1}(T)\leq1+\alpha\). A fresh Gaussian block of width
\(r\) and a polynomial of degree
\(O(\alpha^{-1/2}\log(k/(\alpha\beta)))\) produce an orthonormal
basis \(U\) of width at most \(r\) such that
\begin{equation}\label{eq:strong}
 \norm{(I-UU^\top)T}_2\leq1+3\alpha
\end{equation}
with probability at least \(1-\beta\).
\end{lemma}
\begin{proof}
Let \(B_0=1+\alpha\), \(L_0=1+2\alpha\), and
\[
 p_v(x)=\frac{{\rm Cheb}_v(2x/B_0-1)}
                  {{\rm Cheb}_v(2L_0/B_0-1)}.
\]
On \([0,B_0]\) its magnitude is at most the reciprocal of the
denominator; on \([L_0,2]\) it is at least one. Its denominator is at
least \(\tfrac12\exp(v\sqrt\alpha)\), using
\({\rm Cheb}_v(1+u)=\cosh(v\operatorname{arcosh}(1+u))\) and
\(\operatorname{arcosh}(1+u)\geq\sqrt u\) for \(0\leq u\leq1\).
The algorithm evaluates the scalar denominator by the three-term recurrence,
until it is at least \(2R/\alpha\), where \(R=(100k/\beta)^3\).
Thus it needs only the stated degree and finite arithmetic.

In a fixed eigenbasis of \(T\), let \(H_0\) index eigenvalues above
\(B_0\); its size \(s\) is at most \(r\). For a fresh Gaussian
\(\Omega\in\R^{k\times r}\), with probability at least \(1-\beta\),
\begin{equation}\label{eq:gauss}
 \Omega_{H_0}\Omega_{H_0}^\dagger=I_s,\qquad
 \norm{\Omega_{H_0^c}\Omega_{H_0}^\dagger}_2\leq R.
\end{equation}
An elementary sufficient bound uses the first \(s\) columns of
\(\Omega_{H_0}\). Conditional on the other square-matrix columns,
each column's distance to their span is a standard absolute Gaussian.
A union bound makes all distances at least \(\beta/(Cs)\).
The inverse Frobenius norm is at most \(Cs^{3/2}/\beta\), since its
row norms are reciprocal column distances. Adding columns only improves
the least singular value. Markov's inequality bounds
\(\norm{\Omega_{H_0^c}}_F\) by \(Ck/\sqrt\beta\).
These bounds give \eqref{eq:gauss} with the generous displayed \(R\).
When \(s=0\), no inverse is needed.

Form \(Y_o=p_v(T)\Omega\) by the Chebyshev recurrence. Let \(E_*\)
select the coordinates in \(H_0\) with eigenvalues at least \(L_0\).
For analysis only, right-multiply \(Y_o\) by
\[
 \Omega_{H_0}^\dagger\,
       \operatorname{diag}(p_v(\lambda(T)_{H_0}))^{-1}E_*.
\]
Its head is exactly the strong eigenvector coordinate matrix. Its tail
norm is at most \(\alpha/2\), by \eqref{eq:gauss} and the chosen denominator.
Only selected strong inverse polynomial values enter this norm bound;
all are at most one. Weak head coordinates are exactly zero, so no
inverse of a tiny weak polynomial value is charged.
Consequently the projector \(P_o\) onto \(\range Y_o\) obeys
\(\norm{(I-P_o)U_*}_2\leq\alpha/2\) for the strong invariant space.
The remaining spectrum is below \(L_0\), and \(\norm T_2\leq2\).
The triangle inequality gives
\(\norm{(I-P_o)T}_2\leq L_0+\alpha=1+3\alpha\).
Any exact orthonormal range basis of \(Y_o\) gives \(U\). Rank loss
is harmless. The recurrence may have exponential growth in its degree;
this proof counts exact arithmetic and does not assert numerical stability.
\end{proof}

\section{Threshold search and the physical error transfer}
Maintain the positive bracket \([\ell,u]=[\rho/(4k),4k\rho]\) for
\(\lambda_*\). At \(t=(\ell+u)/2\), call Lemma~\ref{lem:quantile}.
On \emph{below}, set \(\ell=t\); on \emph{above}, set \(u=t\).
These decisions are safe by \eqref{eq:quantile}.
On \emph{near}, \(|\mu_j(C_t)|\leq\alpha\);
stop and use this \(t\). If the bracket reaches
\(u-\ell\leq\alpha h\ell\) first, stop at its upper endpoint \(t=u\).

An uncertain threshold is sufficiently close in a stronger, weighted
sense. Put \(b_* =\alpha/(1-\alpha^2)\leq2\alpha\). For
\(-1\leq c\leq\alpha\),
\(c\leq b_*(1-c^2)\). Apply this inequality on the bottom-\(j\)
eigenspace of \(C_t\), undo \(J\), and use Lemma~\ref{lem:defect}.
On a dimension-\(j\) subspace,
\[
 y^\top K_ty\leq36b_*h\,y^\top Dy.
\]
Since \(D\preceq2(tB_m+Q)\), generalized min--max gives
\begin{equation}\label{eq:near}
 \lambda_*\geq\frac{1-x_*}{1+x_*}t
                \geq(1-288\alpha h)t,
 \qquad x_*=72b_*h\leq144\alpha h.
\end{equation}
The uncertainty also gives \(\lambda_j(T_t)\leq1+\alpha\).
At the final upper endpoint, inertia gives the stronger
\(\lambda_j(T_t)\leq1\), and
\begin{equation}\label{eq:upper}
 \lambda_*\leq t\leq(1+\alpha h)\lambda_*.
\end{equation}
Thus either exit permits Lemma~\ref{lem:filter}.
The initial ratio \(16k^2\) bounds the number of calls by
\(Q_{\max}=\lceil\log_2(16\cdot10^6m^2/\theta)\rceil+1\).
This bound can be found by scalar halving, without a logarithm oracle.
All thresholds are allowed: no pole avoidance or eigenvalue separation
is needed, since the explicit \(D,F\) remain positive definite.

\begin{lemma}[Weighted strong-action transfer]\label{lem:transfer}
For the final \(t\), take \(U\) from Lemma~\ref{lem:filter} with
\(T=T_t\), and set \(\mathcal S=\range(ZJU)\). Then
\begin{equation}\label{eq:transfer}
 \norm{(I-\Pi_{\mathcal S})M}_2^2
 \leq\frac{1+252\alpha h+\varepsilon_m}{1-252\alpha h}\,t.
\end{equation}
\end{lemma}
\begin{proof}
Let \(b=(1+3\alpha)^2-1\leq7\alpha<1\). For \(w\perp\range U\),
\eqref{eq:strong} and symmetry imply
\(w^\top(C_t^2-2C_t)w\leq b\norm w^2\). Substitute \(w=J^\top y\).
The defect identity \(F=K_tF^{-1}K_t+\Delta\) gives
\begin{align*}
 2y^\top K_ty
 &\geq y^\top K_tF^{-1}K_ty-b\,y^\top Fy\\
 &=(1-b)y^\top K_tF^{-1}K_ty-b\,y^\top\Delta y
 \geq-36bh\,y^\top Dy.
\end{align*}
Hence \(y^\top K_ty\geq-126\alpha h\,y^\top Dy\) on this dual
complement. The factor \(b\) multiplying the defect is essential.
Transferring only a lower bound on \(C_t\) would lose the extra \(h\).

For a physical \(x\perp\mathcal S\), put \(y=Z^\top x\). Then
\(J^\top y\perp\range U\), and
\[
 y^\top G^{-1}y\leq\norm x^2,\quad
 \norm y^2\leq2\norm x^2,\quad
 y^\top Q_0y\leq2x^\top MM^\top x.
\]
In particular \(y^\top Dy\leq2t\norm x^2+2x^\top MM^\top x\).
Using the mass upper bound in \eqref{eq:massbounds},
\begin{align*}
 t\norm x^2-x^\top MM^\top x
 &\geq y^\top(tG^{-1}-Q)y\\
 &\geq-(252\alpha h+\varepsilon_m)t\norm x^2
           -252\alpha h\,x^\top MM^\top x.
\end{align*}
Taking the maximum Rayleigh quotient over \(\mathcal S^\perp\)
proves \eqref{eq:transfer}. The correct physical lift is \(ZJU\);
no \(G^{-1}\) or \(J^{-\top}\) occurs in this lift.
\end{proof}

Combine \eqref{eq:transfer}, \eqref{eq:near} or \eqref{eq:upper},
and \eqref{eq:hatcompare}. Their factor over \(\lambda_{r+1}(H)\)
is less than \(1+2000(\alpha h+\varepsilon_m)<1+\theta\), for the
stated parameters. For example the transfer factor is at most
\(1+758\alpha h\), because \(\varepsilon_m=\alpha h\); the uncertain
threshold costs at most \(1+576\alpha h\), and the mass comparison
at most \(1+2\varepsilon_m\). Their product is bounded as stated.
This is stronger than the squared factor \((1+\theta)^2\) needed
for Theorem~\ref{thm:main}.
Orthonormalize \(ZJU\) to obtain \(U_p\), compute \(R_p=M^\top U_p\)
through the original factors, and explicitly output \(C=U_pR_p^\top\).
Its width is at most \(r\), since \(Z\) is injective and \(J\) is
nonsingular on the support.

\section{Costs, confidence, and exact-arithmetic limits}
The sketches and support cost \(\tO(I+G_h+m^{\omega+\zeta})\).
At every threshold, the finite Newton construction and coefficient
extraction cost \(\tO(m^{\omega+\zeta})\). Computing raw Krylov
blocks costs \(\tO(q\Ac(r+1))\). Applying \(T\) to its exact range
basis in width-\((r+1)\) chunks adds the same cost.
The final fresh filter also costs \(\tO(q\Ac(r))\).

One must charge for the \emph{raw} Krylov columns, including dependent
ones. There are at most \(L=(2q+1)(r+1)\), so conventional exact rank
revelation and orthogonalization cost \(O(kL^2)\). The compressed Gram
costs at most that much. This proves the conservative charge \(mL^2\)
in \eqref{eq:general}. If \(L\leq m\), pad the \(k\)-by-\(L\)
matrix to an \(m\)-by-\(m\) matrix, apply deterministic CUP rank
revelation \cite{jps}, and use \eqref{eq:qr} on its independent columns.
This costs \(\tO(m^{\omega+\zeta})\), even with arbitrary dependencies.
The compressed Gram and the final core range basis have the same fast
bound. Gaussian repetitions use separate bases, so only multiply the
cost by logarithms and do not enlarge the raw-width condition.

The physical lift and \(M^\top U_p\) cost \(O(\Ac(r))\).
For its independent physical columns, \eqref{eq:qr} costs
\(\MM(r,n,r)+O(r^{\omega+\zeta})+\MM(n,r,r)\), bounded by
\(O(\Ac(r)+m^{\omega+\zeta})\) by padding and permutation symmetry.
Thus an \(nr^2\) term is unnecessary. Writing and expanding the final
product costs \(np+\MM(n,r,p)\). This proves both runtime bounds.

Allocate \(\delta/2\) to the embeddings, and
\(\beta=\delta/[4(Q_{\max}+1)]\) to each conditional quantile and the
final fresh filter. Each threshold depends only on prior random draws;
the per-query guarantees hold conditionally with fresh draws. Union
bounds prove success at least \(1-\delta\). The cap and all polynomial
degrees depend only on supplied parameters and computed dimensions.
Check ranks and nonzero coefficient denominators before divisions.
An unexpected invalid check returns zero factors within the same cap.
In fact the support \(G,Q_0\), the even inverse-mass polynomial, and
the explicit \(F\) stay positive definite on every computed nonzero
support even off the sketch event; only their approximation bounds can
fail. All routine dimensions satisfy \(\widehat k\leq\rank M\leq m\).

No bit-complexity conclusion follows. Characteristic coefficients can have
large representations, and their signs can require arbitrarily high
precision. The \(D\) inverses, Cholesky, and exact rank tests can also be
arbitrarily ill-conditioned. The final polynomial can grow exponentially in its
degree. Their real operation counts are uniform, but stable finite-word
implementations require a separate theorem.

\section{A growing-rank total-output improvement}
The following consequence uses the current rectangular bound \cite{amanov}.
For sorted exponents \(a\geq b\geq c\), the multiplication exponent
is \(a+b\) when \(c(a+b)\leq ab\), and is at most
\((ab+bc+ca)^2/(4abc)\) otherwise.

\begin{corollary}\label{cor:example}
Let \(n=p=N=m^{59/50+o(1)}\), \(\theta=m^{-21/50}\), and take
constant failure probability. For every
\(1\leq r\leq\lceil m^{59/100}\rceil=N^{1/2+o(1)}\),
Choosing \(h=m^{-7/100}\), Theorem~\ref{thm:main} constructs the
approximation in
\[
 N^{6724/3363+o(1)}=N^{2-2/3363+o(1)}
\]
operations and explicitly outputs it in total \(N^{2+o(1)}\).
The exact-Gram bound is \(N^{7056/3481+o(1)}\).
Throughout the subrange \(r\leq\lceil m^{59/109}\rceil\), the checked
ordinary fast batched block-Krylov bound is \(N^{239/118+o(1)}\).
\end{corollary}
\begin{proof}
In powers of \(m\), \(q=m^{7/40+o(1)}\) and
\(L\leq m^{153/200+o(1)}<m\). At the largest rank, rectangular
multiplication gives \(\Ac(r+1)=m^{43681/20000+o(1)}\); padding
bounds every smaller rank by this cost. The sketch row exponent is
\(1+2(7/100)=57/50<59/50\). The construction exponents are at most
\[
 G_h:\ \frac{3362}{1425},\qquad
 q\Ac(r+1):\ \frac{47181}{20000},\qquad
 m^\omega:\ \frac94.
\]
The first is largest and is less than the output exponent \(59/25\)
by \(1/1425\). It exceeds the action exponent by \(283/1140000>0\).
Dividing by \(59/50\) gives the stated construction exponent. All fixed
multiplication slacks can be chosen inside these positive margins.
Since \(r\leq N^{1/2+o(1)}\), expansion has exponent two.
The fast raw-width core bound is essential at the largest ranks:
conventional \(mL^2\) would have exponent \(253/100\), above output.
This is an exponent-two statement, not a literal \(O(N^2)\) assertion.

In the stated smaller-rank subrange, each ordinary block-Krylov factor
action is input optimal, of exponent \(109/50\), and its iteration count has
exponent \(21/100\). Its checked bound is therefore
\(m^{239/100+o(1)}=N^{239/118+o(1)}\). At its upper endpoint the
raw basis has exponent \(8189/10900\). Fast tall rank revelation costs
at most \(m^{92393/43600+o(1)}\), and the compressed Gram and even cubic
Ritz work have exponents less than \(239/100\). Wider factor actions
and projection also lie below that bound. At larger ranks, use the
exact-Gram comparison \(N^{7056/3481+o(1)}\); the ordinary action and
Ritz costs can change. For example, at \(r=\lceil m^{1/10}\rceil\),
the previously developed scalar stochastic route has exponent \(489/200\)
in powers of \(m\), also above output. These are comparisons of
available upper bounds, not lower bounds on all algorithms.
\end{proof}

\section{Absolute error without a supplied target rank}
When the desired error is absolute, the threshold is known and no
coefficient-ratio initialization or threshold bisection is needed.

\begin{corollary}[Adaptive absolute error]\label{cor:absolute}
Given \(e>0\) and the parameters of Theorem~\ref{thm:main}, define
\begin{gather*}
 r_e=\min\{s\in\{0,\ldots,m\}:\sigma_{s+1}(M)\leq e/(1+\theta)\},\\
 \ell_e=\min(m,2r_e),\quad s_e=\min(m,2r_e+1).
\end{gather*}
With probability at least \(1-\delta\), an algorithm, without being
supplied \(r_e\), explicitly returns an approximation of error at most
\(e\) and width at most \(2r_e\). On that same event, if
\((2q+1)s_e\leq m\), its total cost is
\begin{equation}\label{eq:absolute}
 \tO\!\left(I+G_h+m^{\omega+\zeta}+q\Ac(s_e)
                  +np+\MM(n,\ell_e,p)\right).
\end{equation}
Every failed outcome has a dimension-and-parameter-dependent arithmetic
cap, using the general raw-width charge when necessary. The
\(r_e\)-sensitive cost above is a high-probability bound.
\end{corollary}
\begin{proof}
Fix \(t=e^2/(1+\theta/4)\), construct its \(F,J,T\) once, and
try target widths \(r=0,1,2,4,\ldots\). Use the exact-support branch
as soon as \(r\geq\widehat k\). At any earlier width, form independent
Musco--Musco spaces with target \(j=r+1\), as in
Lemma~\ref{lem:quantile}. If any compressed \(\beta_j>1\), advance
the rank. Otherwise, when any repetition has the per-vector guarantee,
\(\lambda_j(T)^2\leq(1-\alpha/100)^{-1}\), so
\(\lambda_j(T)\leq1+\alpha\). Apply Lemma~\ref{lem:filter} at
width \(r\), and use the same dual physical lift. The factor in
\eqref{eq:transfer} is at most \(1+758\theta/10^6<1+\theta/4\),
so the squared error is at most \(e^2\).
At \(r=0\), no projection filter is needed: the accepted bound is
\(\norm T_2\leq1+\alpha\), and the zero-dimensional space already
satisfies the strong residual certificate. Write zero explicitly.

By \eqref{eq:hatcompare} and the definition of \(r_e\),
\[
 \lambda_{r_e+1}(\widehat H)
 \leq\frac{e^2}{(1+\theta)^2(1-\varepsilon_m)}<t.
\]
Consequently at every attempted \(r\geq r_e\) with \(r<k\),
\(\lambda_{r+1}(T)\leq1\). Compression min--max then gives
\(\beta_{r+1}\leq1\) on every repetition, independently of its quality.
The rank loop therefore stops by its first width at least \(r_e\),
which is at most \(2r_e\); if \(r_e=0\), it stops at zero.
An intervening exact-support branch only reduces the returned width.
This stopping-rank bound needs only the common sketch event.

There are at most \(\lceil\log_2m\rceil+2\) attempts. Allocate fresh
conditional per-vector and final-filter failure budgets across them.
Their union with the sketches has failure probability at most \(\delta\).
Every processed block width is at most \(s_e\) on that event, so
monotonicity and the logarithmic number of attempts give
\eqref{eq:absolute}. The displayed raw-width condition justifies fast
range algebra throughout. Off the sketch event, stop the rank loop at
\(\widehat k\leq m\) and use the general charge from
Theorem~\ref{thm:main}; no unbounded retries or precision loops occur.
\end{proof}

\section{Antecedents and scope of the claim}
Absolute-value SPD preconditioning is due to Vecharynski--Knyazev
\cite{vk}. Vecharynski--Yang \cite{vy} already combine inertia-preserving
congruence, AV preconditioning, and randomized Lanczos or Arnoldi
eigenvalue counts. Block PLHR \cite{plhr} explicitly uses the ideal
squared-sign nullspace identity. We do not claim these ideas as new.
Their inspected theorems do not provide this complete arbitrary-factor,
gap-free, width-\(r\), relative spectral-tail runtime. In particular an
AV spectral-equivalence theorem with a positive separation from zero
cannot apply uniformly at our singular thresholds. Our \(F\) is positive
even when \(K_t\) is singular.

The quantitative obligation here is \(0\preceq F-K_tF^{-1}K_t\preceq36hD\)
and its transfer with the additional factor \(\alpha\), followed by
the physical lift \(ZJU\). Bounded congruent spectrum and inertia alone
do not imply this: \(F=D=I\), \(K=-\alpha I\) has both properties,
but no \(O(\alpha h)D\) lower bound for arbitrarily small \(h\).
The per-vector theorem used for ordered quantiles is exactly the
published Musco--Musco theorem. Gaussian range finding, Chebyshev filters,
Descartes' rule, and fast characteristic-polynomial algorithms are also
established tools. The inspected sources do not establish the combined
runtime, but that audit is not a global priority certificate.

The original harder target with \(m=N^{8/9}\) and \(\theta=1/m\)
remains open. This manuscript proves a different unrestricted-input
parameter regime; it does not silently replace that target or assert
a stable numerical implementation.

\begin{thebibliography}{99}
\bibitem{amanov} A. Amanov, M. van den Berg, J. Zuiddam.
\emph{Multiplying matrices via polynomials: exposition and barriers}. 2026.
\url{https://arxiv.org/abs/2610.30054}.
\bibitem{cnw} M. Cohen, J. Nelson, D. Woodruff.
\emph{Optimal Approximate Matrix Product in Terms of Stable Rank}. 2016.
\url{https://arxiv.org/abs/1507.02268}.
\bibitem{jps} C.-P. Jeannerod, C. Pernet, A. Storjohann.
\emph{Rank-profile revealing Gaussian elimination and the CUP matrix
decomposition}. J. Symbolic Computation, 2013. Algorithm 6, Propositions 3, 6.
\url{https://arxiv.org/abs/1112.5717}.
\bibitem{vk} E. Vecharynski, A. V. Knyazev.
\emph{Absolute value preconditioning for symmetric indefinite linear systems}.
SIAM J. Sci. Comput. 35(2), A696--A718, 2013.
\url{https://arxiv.org/abs/1104.4530}.
\bibitem{vy} E. Vecharynski, C. Yang.
\emph{Preconditioned iterative methods for eigenvalue counts}. 2016/2017.
\url{https://arxiv.org/abs/1602.02306}.
\bibitem{plhr} E. Vecharynski, A. Knyazev.
\emph{Preconditioned Locally Harmonic Residual Method for Computing Interior
Eigenpairs of Certain Classes of Hermitian Matrices}.
SIAM J. Sci. Comput. 37(5), S3--S29, 2015.
\url{https://arxiv.org/abs/1408.0042}.
\bibitem{mm} C. Musco, C. Musco.
\emph{Randomized Block Krylov Methods for Stronger and Faster Approximate
Singular Value Decomposition}. NeurIPS 2015. Algorithm 2, Theorem 12.
\url{https://arxiv.org/abs/1504.05477}.
\bibitem{np} V. Neiger, C. Pernet.
\emph{Deterministic computation of the characteristic polynomial in the
time of matrix multiplication}. J. Complexity, 2021. Theorem 1.1.
\url{https://arxiv.org/abs/2010.04662}.
\bibitem{npv} V. Neiger, C. Pernet, G. Villard.
\emph{Computing Krylov iterates in the time of matrix multiplication}.
ISSAC 2024. Section 5.1, Corollary 5.2.
\url{https://arxiv.org/abs/2402.07345}.
\end{thebibliography}
\end{document}
